\NeedsTeXFormat{LaTeX2e}

\documentclass[portrait]{a0poster}

\usepackage{epsfig}
\usepackage{color}
\usepackage{transparent}
\usepackage{wallpaper}
\usepackage{amsmath}
\usepackage{amssymb}
\usepackage{calc}
\usepackage{bm}
\usepackage[utf8]{inputenc}
\usepackage{textcomp}
\usepackage{mathptmx}
\usepackage[scaled=.90]{helvet}
\usepackage{courier}
\usepackage{listings}

\hyphenation{re-so-na-tors na-no-tubes}

\pagestyle{empty}
\TileWallPaper{10.234cm}{11.594cm}{images/dokutitle}

\setlength{\parskip}{5pt plus 2pt minus 1pt}
\setlength{\parindent}{0pt}

\frenchspacing
\sloppy

\renewcommand{\familydefault}{\sfdefault}

% some color definitions
\definecolor{darkblue}{rgb}{0, 0, 0.5}
\definecolor{lightblue}{rgb}{0, 0, 0.8}
\definecolor{lightgray}{gray}{0.75}

\newcommand{\blue}{\color[rgb]{0.2,0.2,1}}
\newcommand{\lightblue}{\color{lightblue}}
\newcommand{\red}{\color{red}}
\newcommand{\darkgreen}{\color[rgb]{0,0.65,0}}

% Delft blauw 
\definecolor{background}{rgb}{1,1,1}
\definecolor{text}{rgb}{0, 0, 0.3}
\definecolor{heading}{rgb}{0, 0, 0.8}

\definecolor{invheading}{rgb}{0.8,0.8,1}

\pagecolor{black}

\lstset{
  showstringspaces=false,
  basicstyle=\ttfamily,
  identifierstyle=\color[rgb]{0,0,0.5},
  stringstyle=\color[rgb]{0.5,0,0.5}\textbf,
  commentstyle=\color[rgb]{0.9,0,0}\textbf
}




% large heading in text
\newcommand{\heading}[1]{
  {\color{heading}\boldmath\textbf{\huge #1}}\\[\medskipamount]
}

% small heading in text
\newcommand{\smallheading}[1]{
  {\color{heading}\textbf{\Large #1}}\\
}

% very small heading in text
\newcommand{\verysmallheading}[1]{
  {\color{heading}\textbf{\large #1}}\\
}

% white bar across the column
\newcommand{\whitebar}{%
\vspace*{0.5cm}
\hspace*{-1.5cm}
\fcolorbox{white}{white}{\rule{30cm}{-0.7cm}}

\vspace*{1cm}
}

% white background for plots etc.
\newcommand{\whitebg}[1]{
 {
  \setlength{\fboxsep}{5mm}
  \setlength{\fboxrule}{1pt}
  \fcolorbox{black}{white}{\color{black}{#1}}
 }
}

% place a figure
\newcommand{\placefigure}[1]{ 
 {
  \vspace{-0.4cm}
  \begin{center}
  \hspace*{-1cm}\begin{minipage}{0.96\textwidth}
   \whitebg{
     \epsfig{#1,width=0.96\textwidth}
   }
   \vspace*{-1cm}
   {\\}
  \end{minipage}
  \end{center} 
  \vspace{-1cm}
 }
}

% place a listing
\newcommand{\placelisting}[1]{ 
 {
  \vspace{-0.2cm}
  \begin{center}
  \hspace*{-1cm}\begin{minipage}{0.96\textwidth}
   \footnotesize \renewcommand{\baselinestretch}{0.85}
   \whitebg{
   \parbox{0.97\textwidth}{
    \lstinputlisting[language=Perl]{#1}
   }
   }
   \renewcommand{\baselinestretch}{1}
   \vspace*{-1cm}
   {\\}
  \end{minipage}
  \end{center} 
  \vspace{-1cm}
 }
}

% place a listing
\newcommand{\placesmalllisting}[1]{ 
 {
  \vspace{-0.2cm}
  \begin{center}
  \hspace*{-1cm}\begin{minipage}{0.96\textwidth}
   \scriptsize \renewcommand{\baselinestretch}{0.8}
   \whitebg{
   \parbox{0.97\textwidth}{
    \lstinputlisting[language=Perl]{#1}
   }
   }
   \renewcommand{\baselinestretch}{1}
   \vspace*{-1cm}
   {\\}
  \end{minipage}
  \end{center} 
  \vspace{-1cm}
 }
}

% place a text
\newcommand{\placetext}[1]{ 
 {
  \vspace{-0.2cm}
  \begin{center}
  \hspace*{-1cm}\begin{minipage}{0.96\textwidth}
   \footnotesize \renewcommand{\baselinestretch}{0.85}
   \whitebg{
   \parbox{0.97\textwidth}{
    \lstinputlisting{#1}
   }
   }
   \renewcommand{\baselinestretch}{1}
   \vspace*{-1cm}
   {\\}
  \end{minipage}
  \end{center} 
 }
}




\begin{document}

\color{text}
\Large
%
\newlength{\seplength}
\newlength{\headerheight}
\newlength{\columnheight}
\newlength{\columnheighta}
\newlength{\columnheightb}
%
\setlength{\columnheight}{107cm}
\setlength{\columnheighta}{\columnheight-6cm}
\setlength{\columnheightb}{\columnheight+6.5cm}
\setlength{\columnwidth}{0.31\textwidth}
%
\setlength{\fboxsep}{10mm}
\setlength{\fboxrule}{0mm}
%
%
%
% the page header
%
\hspace*{-1cm}\begin{minipage}[t][][t]{\textwidth-2\fboxsep-8\fboxrule}
\color{invheading}
\begin{center}
{
\VERYHuge \vspace*{-0.6cm}
\textsf{\textbf{
Lab::Measurement -- measurement control with Perl
}}
}
\\[\baselineskip]
\begin{tabular}{ccc}
\parbox{5cm}{\vspace*{-1cm}
\begin{center}
%\epsfig{file=images/logo-ur,height=4.5cm}
\end{center}
}
&
\parbox{70cm}{
\begin{center} 
{
\vspace*{-0.5cm}
\huge 
M. Schambeck, E. Fabrizzi, F. Weinelt,  S. Reinhardt, and A. K. Hüttel
}
\\[\medskipamount]
{\it 
\Large
Institute for Experimental and Applied
Physics, University of Regensburg, 93040 Regensburg, Germany\\
}
\end{center}
\vspace*{0.5cm}
}
&
\parbox{5cm}{\vspace*{-1cm}
\begin{center}
%\epsfig{file=images/logo-sfb689,height=4.5cm}
\end{center}
}
\end{tabular}
\end{center}
\end{minipage}
\vspace*{9mm}





%
% begin of first row first column
\fcolorbox{white}{background}{
\begin{minipage}[t][\columnheighta-2\fboxsep-2\fboxrule][t]
  {\columnwidth-2\fboxsep-2\fboxrule} \rule{0pt}{0pt}\\
  \begin{minipage}{\textwidth}
%
%

\vspace*{0.9cm}
\heading{Flexible measurement needed?!}

\vspace*{-1cm}
\begin{itemize}
\item
Tired of following your wires in square meters of LabVIEW diagrams?
\item
Tired of clumsy string handling and low-level driver functions
in your looong C program?
\item
Use a text processing language to manage your measurement! Use Perl!
\end{itemize}
\vspace*{0.5cm}

\placelisting{srs_read.pl}

\vspace*{2cm}
\vspace*{0.9cm}
\heading{Currently supported hardware}

\vspace*{-3mm}
\placefigure{file=images/hardware}
Hardware driver backends: 
\begin{itemize}
\item 
  NI-VISA (MS Windows) and all hardware supported by it
\item
  LinuxGPIB and all hardware supported by it
\item
  TCP connection, generic network socket
\item
  USB-TMC lightweight driver (Linux, libusb)
\item
  VXI-11 lightweight driver (Linux, libtirpc)
\item
  Oxford Instruments IsoBus
\item
  Zurich Instruments LabOne API
\end{itemize}

\vspace*{1cm}
Growing number of high-level drivers (more are {\it very} easy to add):
{\large
\begin{itemize}
\item Multimeters: HP / Agilent / Keysight
\item DC sources: Yokogawa / Keithley / Keysight
\item AW generators: Rigol
\item Lock-in amplifiers: Stanford Research / Signal Recovery / Zurich 
Instruments / SyncTek
\item Temperature controllers: Lakeshore / Oxford Instruments
\item RF / microwave sources, spectrum analyzers, VNAs: \\ Rohde \& Schwarz / 
HP / Agilent / Rigol
\item Nanonis Tramea
\item and many more...
\end{itemize}
}
\vspace*{1cm}
\heading{Key facts}
\vspace*{-2cm}
\begin{itemize}
 \item
 Open source / free software\\[-5cm]
 \item 
 {\lightblue https://www.labmeasurement.de/}
 \hfill\epsfig{file=images/qrcode-labmeasurement,width=5cm}
 \item
 License: same as Perl (GPL-1+ or Artistic)
 \item
 Releases on CPAN, development on Github
 \item
 Contributors and cooperations welcome!
\end{itemize}

%
%
%
\end{minipage}
\end{minipage}}
% end of first column
%
%
% begin of second column
\fcolorbox{white}{background}{
\begin{minipage}[t][\columnheighta-2\fboxsep-2\fboxrule][t]
  {\columnwidth-2\fboxsep-2\fboxrule} \rule{0pt}{0pt}\\
  \begin{minipage}{\textwidth}
%
%
%

\vspace*{0.9cm}
\heading{Real world measurement}

\vspace*{-1cm}
\begin{itemize}
 \item 
Ferromagnetic resonance measurement with vector network analyzer (VNA)
 \item
Outer loop: continuous sweep of magnetic field with OI Mercury iPS magnet 
controller
 \item 
Inner loop: VNA transmission measurement at multiple discrete microwave 
frequencies
\end{itemize}

\vspace*{0.5cm}
 \placesmalllisting{VNA-FMR.pl}
\vspace*{3cm}

\heading{Output files}
\vspace*{-1cm}
\placelisting{ls.txt}
\vspace*{2cm}
\begin{itemize}
\item \emph{META.yml:} various metadata (host, user, date, L::M version,
  command line arguments)
\item \emph{data.dat:} measured data, in tab-separated Gnuplot format
\item \emph{data.png:} live plot at the end of the measurement, as a png image
\item \emph{magnetic-resonance.pl:} archival copy of the measurement script

\end{itemize}
\vspace*{2cm}
\heading{Advanced sweep features}

\vspace*{-1cm}
\begin{itemize}
\item Multidimensional sweeps, e.g. 3D sweep: creating one 2D datafile for each
  step of the outermost sweep
\item Log arrays and matrices of data (PDLs). Useful for spectrum analyzers,
  VNAs, oscilloscopes which do fast, internally controlled sweeps.
\item Extensive support for live plots via gnuplot: line plots (2D data) and 
  color maps (3D data)
\item Customizing live plots: access to all gnuplot plot and curve options 
  via PDL::Graphics::Gnuplot
\end{itemize}


% 
% 
% 
\end{minipage}
\end{minipage}}
% end of second column
%
%
%
%begin third column
\fcolorbox{white}{background}{
\begin{minipage}[t][\columnheighta-2\fboxsep-2\fboxrule][t]
  {\columnwidth-2\fboxsep-2\fboxrule} \rule{0pt}{0pt}\\
  \begin{minipage}{\textwidth}
%
%
%

\vspace*{0.9cm}
\heading{Layer structure}

\vspace*{-3mm}
\placefigure{file=images/structure}

\begin{itemize}
\item
Modular structure. Easy to extend with new instrument drivers and connection
types
\item
Abstract IO layer, makes instrument drivers independent of hardware backends
\end{itemize}

\vspace*{2cm}
\heading{High-level sweep framework}
\vspace*{-1cm}

\begin{itemize}
\item Modern Perl implementation; use state of the art object-oriented
  programming
  \vspace*{1cm}
  \placefigure{file=images/moose-sweep}
  \vspace*{0.5cm}
\item Separate classes for sweeps, datafiles, datafolders, and plots
\item Most operational details of sweeps implemented in subclasses of
  Lab::Moose::Sweep
\item High modularity: \emph{very} easy to extend
\end{itemize}
\vspace*{2cm}
  \heading{New improvements / additions}
  \vspace*{-2cm}
  \begin{itemize}
  \item
  VISA::USB connection: USB-TMC over NI-VISA
  \item
  Enable magnet persistent mode usage
  \item
  RFC3161 timestamps for measurement data
  \item {\lightblue New drivers}
  \begin{itemize}
  \item 
  Andeen-Hagerling AH2700A ultra-precision capacitance bridge
  \item
  HP 8596E, Rigol DSA815 spectrum analyzer
  \item
  SignalRecovery 7265, Zurich Instruments HF2LI lock-in
  \item
  Oxford Instruments Triton and ITC503 temperature control
  \item
  Keithley 2400, 2450 sourcemeter
  \item
  HP 34420A Nanovoltmeter
  \item
  Tektronix TBS1000C, Keysight DSOS604A oscilloscope
  \item
  Rigol DG5000 series waveform generator
  \item
  Rohde \& Schwarz ZNL network analyzer
  \item
  Nanonis Tramea (work in progress)
  \end{itemize}
\end{itemize}

\vspace*{2cm}
\heading{Reference / Cite as}
\textbf{``Lab::Measurement --- a portable and extensible framework for 
controlling lab  equipment and conducting measurements''}, 
S. Reinhardt {\it et al.}, Comp. Phys. Comm. \textbf{234}, 216 (2019)

\end{minipage}
\end{minipage}}  % end of first row third column
%
%
%

\vspace*{0.5cm}
\begin{center}\color{invheading}\large
\hspace*{-4cm}We gratefully acknowlegde funding by the DFG via the Emmy Noether 
grant Hu1808/1, SFB 689, SFB 1277, and GRK 1570.
\end{center}

\end{document}
